\section{Related Work}

\textbf{LLM-driven fuzzing and test generation.} LLMs have been used as
zero-shot input generators for deep-learning libraries~\cite{deng2023titanfuzz},
as universal fuzzers across languages~\cite{xia2024fuzz4all}, and to escape
coverage plateaus in search-based test generation~\cite{lemieux2023codamosa}.
These systems draw on what the model already knows from pre-training. Our
earlier studies put such a model inside a bounded, coverage-free refinement
loop, first against C parsers~\grammarcite{} and then against Dart
deserializers with a differential oracle~\dartcite{}; the second is where
the knowledge gap studied here was measured.

\textbf{Knowledge from reading code.} WhiteFox~\cite{yang2024whitefox} has
an LLM agent read compiler optimization source code and write requirements
for test programs that trigger it. KernelGPT~\cite{yang2025kernelgpt}
infers syscall specifications from kernel source for Syzkaller.
DiffSpec~\cite{rao2024diffspec} generates differential tests from natural-language
specifications and code artifacts. AIJon~\cite{vadayath2026aijon} generates
IJON annotations from source and finds them roughly on par with human-written
ones in a coverage-guided fuzzer. All of these derive knowledge by reading
code or specifications, for coverage or for test requirements. None compares
knowledge sources under control, and none acquires knowledge by running the
target. Our \emph{code} arm is a small, controlled instance of this family;
our \emph{probe} arm is the black-box alternative.

\textbf{Differential testing of JSON processing.} Differential
testing~\cite{mckeeman1998differential} of JSON parsers across languages has
found many semantic discrepancies at the parsing level, text to
tree~\cite{moller2024crosslanguage}. JsonATG~\cite{wang2025jsonatg} uses a
multi-agent LLM pipeline to generate bug-triggering tests for Java JSON
libraries. We test \emph{typed binding}, tree to class, where divergences
come from how each library maps JSON onto a language's types and
null-safety rules, and we compare four libraries on one class in two
ecosystems.

\textbf{Known Kotlin behaviour.} That Gson bypasses Kotlin null-safety is a
long-standing known issue~\cite{gsonissue1657}, which we use as an answer
key: a behaviour the agents should rediscover without being told.
